% theme: changes_of_state
\MajorQuestion{物質の状態変化}
\SetRowSpaceQanda

\Instruction{物質の状態変化と質量・体積に関する次の問いに答えなさい。}
\QQ{温度によって物質の状態が変化することを\NamedBlank{state_change}という。}{}
\QQ{固体・液体・気体の3つの状態をまとめて何というか。}{}
\QQ{固体から直接気体に変化したり、気体から直接固体に変化したりする\RefBlank{state_change}を何というか。}{}
\QQ{物質が\RefBlank{state_change}するとき、変化しない量は質量と体積のどちらか。}{}
\QQ{一般に、固体・液体・気体のうち体積が最も大きい状態はどれか。}{}

\Instruction{物質の状態と粒子のようすに関する次の問いに答えなさい。}
\QQ{物質をつくっている非常に小さなものを何というか。}{}
\QQ{粒子が規則正しく並び、その場で振動している状態を\NamedBlank{solid}という。}{}
\QQ{粒子が近くに集まり、互いに動き回っている状態を\NamedBlank{liquid}という。}{}
\QQ{粒子が離れて存在し、自由に飛び回っている状態を\NamedBlank{gas}という。}{}
\QQ{水が\RefBlank{solid}の氷になるとき、体積は大きくなるか小さくなるか。}{}

\Instruction{融点、沸点、蒸発と沸騰に関する次の問いに答えなさい。}
\QQ{\RefBlank{solid}がとけて\RefBlank{liquid}に変化するときの温度を\NamedBlank{melting_point}という。}{}
\QQ{\RefBlank{liquid}が沸騰して\RefBlank{gas}に変化するときの温度を\NamedBlank{boiling_point}という。}{}
\QQ{\RefBlank{liquid}が液面から\RefBlank{gas}に変化することを何というか。}{}
\QQ{\RefBlank{liquid}の内部からも\RefBlank{gas}に変化することを何というか。}{}
\QQ{一定の\RefBlank{melting_point}や\RefBlank{boiling_point}を示すのは、純粋な物質と混合物のどちらか。}{}

\Instruction{混合物の温度変化と蒸留に関する次の問いに答えなさい。}
\QQ{一定の融点や沸点を示さず、状態変化中も温度が変化するのは何か。}{}
\QQ{液体を沸騰させ、出てきた気体を冷やして再び液体として取り出す方法を\NamedBlank{distillation}という。}{}
\QQ{\RefBlank{distillation}では、液体どうしの何の違いを利用して物質を分けるか。}{}
\QQ{沸点が約78℃で、水と混ざり合う液体を\NamedBlank{ethanol}という。}{}
\QQ{水と\RefBlank{ethanol}の混合物を蒸留したとき、はじめに得られる液体に多く含まれる物質は何か。}{}
